% 核心结论：电场和磁场的旋度分别联系穿面磁场变化和总电流。
% 图契约：物理竞赛讲义；印刷宽16cm，最小字号8.5pt；几何示意与公式分区。
% 布局：[几何/曲线 7cm] [定义/结果 7cm]；实线主量，虚线辅助；无装饰数据。
\documentclass[tikz,border=6pt]{standalone}
\usepackage{fontspec,xeCJK,amsmath,amssymb,graphicx}
\IfFontExistsTF{Times New Roman}{\setmainfont{Times New Roman}}{\setmainfont{TeX Gyre Termes}}
\IfFontExistsTF{Heiti SC}{\setCJKmainfont[AutoFakeBold=1.5]{Heiti SC}}{\setCJKmainfont[AutoFakeBold=1.5]{FandolHei-Regular}}
\usetikzlibrary{arrows.meta,bending,calc,positioning,patterns}
\definecolor{ink}{HTML}{243745}
\definecolor{blue}{HTML}{275DAB}
\definecolor{green}{HTML}{177C63}
\definecolor{red}{HTML}{BD4935}
\definecolor{orange}{HTML}{AA6C16}
\definecolor{purple}{HTML}{715393}
\tikzset{
 every node/.style={font=\fontsize{9.5}{13}\selectfont,text=ink},
 axis/.style={-{Stealth[length=4pt]},line width=.65pt,draw=ink!65},
 vec/.style={-{Stealth[length=5pt]},line width=1pt,draw=blue},
 short/.style={-{Stealth[length=3pt]},line width=.8pt},
 guide/.style={dashed,draw=ink!40,line width=.5pt},
 card/.style={rounded corners=2pt,draw=blue!60,fill=blue!4,
   align=center,text width=6.8cm,inner xsep=9pt,inner ysep=7pt},
 title/.style={font=\fontsize{11}{14}\selectfont\bfseries,text width=15cm,align=center},
 label/.style={font=\fontsize{8.5}{11}\selectfont,align=center}}
\begin{document}
\begin{tikzpicture}[x=1cm,y=1cm]
\path[use as bounding box] (-8,-3.35) rectangle (8,3.35);
\node[title] at (0,2.95) {麦克斯韦旋度方程：环量与穿面变化};
\begin{scope}[shift={(-4,0)}]

\draw[blue!30,fill=blue!3] (0,0) circle (1.8);
\foreach \a in {0,45,...,315}{\draw[short,blue] (\a:1.8) arc[start angle=\a,end angle=\a+20,radius=1.8];}
\node[blue] at (0,0) {$\odot\,\mathbf e_z$};
\node[label] at (0,-2.45) {边界正向与曲面法向按右手定则绑定。};
\end{scope}\node[card,draw=blue!65,fill=blue!4] at (3.9,1.65) {$\displaystyle \nabla\times\mathbf E=-\frac{\partial\mathbf B}{\partial t}$\\变化磁场产生有旋电场。};
\node[card,draw=green!65,fill=green!4] at (3.9,0) {$\displaystyle \nabla\times\mathbf B=\mu_0\mathbf J+\mu_0\varepsilon_0\frac{\partial\mathbf E}{\partial t}$\\传导电流与位移电流共同贡献。};
\node[card,draw=orange!65,fill=orange!4] at (3.9,-1.65) {固定曲面上的积分式经斯托克斯定理局域化；先检查场在邻域内的光滑性。};
\end{tikzpicture}\end{document}
